\documentclass{report}
\usepackage{amsmath,amssymb}
\setlength{\parindent}{0mm}
\setlength{\parskip}{1em}
\begin{document}
\begin{center}
\rule{6in}{1pt} \
{\large Luis Chacon \\
{\bf An optimal asymptotic-preserving solver for the extremely anisotropic heat transport equation in magnetized plasmas}}

Oak Ridge National Laboratory \\ PO Box 2008 \\ Oak Ridge \\ TN 37831
\\
{\tt chaconl@ornl.gov}\\
Diego del-Castillo-Negrete\end{center}

Transport in magnetized plasmas is of fundamental interest in controlled
fusion and astrophysics research. Two issues make this problem particularly
difficult to study: (i) The \emph{extreme anisotropy} between the
parallel (i.e., along the magnetic field), $\chi_{\parallel}$, and
the perpendicular, $\chi_{\perp}$, conductivities ($\chi_{\parallel}/\chi_{\perp}$
may exceed $10^{10}$ in fusion plasmas); and (ii) \emph{magnetic}
\emph{field-line chaos,} which in general precludes the construction
of magnetic field line coordinates. In fact, to date, and despite
significant work on the subject (see e.g. {[}1-6{]}), a suitable numerical
approach that is at the same time high-order (to avoid numerical pollution
of the perpendicular dynamics), robustly positivity-preserving (i.e.,
that enforces a maximum principle and thus a positive temperature
at all times), and algorithmically scalable (i.e., amenable to modern
iterative methods) does not exist. Approaches that enforce the maximum
principle of the parallel transport operator in finite differences
{[}5{]} and finite elements {[}6{]} have been proposed. However, they
rely on limiters which result in formally first-order accurate discrete
representations (thus introducing numerical pollution) and result
in strongly nonlinear algebraic systems which tend to break modern
nonlinear iterative methods. Moreover, the condition number of the
Jacobian matrix associated with the anisotropic transport equation
can be shown to scale as the anisotropy ratio $\chi_{\parallel}/\chi_{\perp}\ggg1$,
thus effectively precluding the use of scalable iterative methods.

Recently {[}7{]}, a novel Lagrangian Green\textquoteright{}s function
method has been proposed to solve the \emph{purely}
parallel transport equation, which is applicable to both integrable
and chaotic magnetic fields. The approach excels in all counts, namely,
it is very accurate (in fact, it respects transport barriers --flux
surfaces-- exactly by construction), inherently positivity-preserving,
and scalable algorithmically (i.e., work per degree-of-freedom is
grid-independent). However, it is of limited applicability to practical
applications, as it does not account for more general physics such
as perpendicular transport and sources.

In this talk, we will review the Lagrangian strategy for purely parallel
transport, and describe its extension to include perpendicular transport
and arbitrary sources. The formulation is asymptotic-preserving (AP)
by construction, which ensures a consistent numerical discretization
temporally and spatially for \emph{arbitrary} $\chi_{\parallel}/\chi_{\perp}$
ratios (from arbitrarily large to $\mathcal{O}(1)$ values). This
is of particular importance, as parallel and perpendicular transport
terms may become comparable in particular regions of the plasma (e.g.,
at incipient islands), while remaining disparate elsewhere. Algorithmically,
the approach uses a simple operator-split approach, consisting of
an Eulerian update (employing multilevel-preconditioned Newton-Krylov
methods), and a Lagrangian step. Both steps scale optimally in that
the cost per degree-of-freedom is grid-independent. The AP character
of the formulation ensures that the splitting error is small in all
regimes, and in fact it becomes negligible in the asymptotic regime
$\chi_{\parallel}/\chi_{\perp}\ggg1$. We will present numerical evidence
that demonstrates the advertised properties of the approach.

\medskip{}


\noindent {[}1{]} C. Sovinec et al., \emph{J. Comput. Phys.}, \textbf{195}
(2004) 355\textendash{}386

\noindent {[}2{]} S. G\"unter, Q. Yu, J. Kr\"uger, K. Lackner, \emph{J.
Comput. Phys.}, \textbf{209} (2005) 354\textendash{}370

\noindent {[}3{]} S. G\"unter, K. Lackner, C. Tichmann, \emph{J. Comp.
Phys.}, \textbf{226} (2), 2306-2316 (2007)

\noindent {[}4{]} S. G\"unter, K. Lackner, \emph{J. Comp. Phys.}, \textbf{228}
(2), 282-293 (2009)

\noindent {[}5{]} P. Sharma, G. W. Hammett, \emph{J. Comput. Phys.},
\textbf{227} (2007) 123\textendash{}142

\noindent {[}6{]} D. Kuzmin, M.J. Shashkov, D. Svyatskiy, \emph{J.
Comput. Phys.}, \textbf{228} (2009) 3448\textendash{}3463

\noindent {[}7{]} D. del-Castillo-Negrete, L. Chac�n, \emph{Phys.
Rev. Lett.}, 195004 (2011)


\end{document}
